\documentclass[11pt]{article}
% ===========================================================================
%  Shared preamble for the AP Calculus AB exam-preparation worksheets.
%  Input right after \documentclass. Files beginning with "_" are not compiled
%  on their own. Navy and gold, to match the course's lesson pages.
% ===========================================================================
\usepackage[letterpaper,top=2.3cm,bottom=2.2cm,left=1.7cm,right=1.7cm,headheight=28pt,footskip=30pt]{geometry}
\usepackage{amsmath,amssymb}
\usepackage[T1]{fontenc}
\usepackage[utf8]{inputenc}
\usepackage{lmodern}
\usepackage{xcolor}
\usepackage[most]{tcolorbox}
\usepackage{enumitem}
\usepackage{fancyhdr}
\usepackage{array}
\usepackage{tikz}
\usepackage{pgfplots}
\pgfplotsset{compat=1.18}

\definecolor{navy}{HTML}{161E54}
\definecolor{ink}{HTML}{27357E}
\definecolor{gold}{HTML}{B98A00}
\definecolor{goldbg}{HTML}{FFFDF3}
\definecolor{inkbg}{HTML}{F6F7FD}
\definecolor{rule}{HTML}{C9CFF0}
\definecolor{okgreen}{HTML}{1B7A44}
\definecolor{okbg}{HTML}{F1F8F3}

\pagestyle{fancy}
\fancyhf{}
\fancyhead[L]{\footnotesize Name: \underline{\hspace{4cm}}}
\fancyhead[C]{\textbf{\color{navy}Integration Academy}\\[-2pt]{\scriptsize Calculus AB \textperiodcentered{} AP\textsuperscript{\textregistered} Exam Preparation}}
\fancyfoot[L]{\tiny\color{black!55}\parbox[t]{0.78\textwidth}{AP\textsuperscript{\textregistered} is a registered trademark of the College Board, which was not involved in the production of, and does not endorse, this material. \textcopyright{} Integration Academy.}}
\fancyfoot[R]{\scriptsize\thepage}
\renewcommand{\headrulewidth}{1.2pt}
\renewcommand{\headrule}{\hbox to\headwidth{\color{gold}\leaders\hrule height \headrulewidth\hfill}}
\renewcommand{\footrulewidth}{0pt}
\setlength{\parindent}{0pt}
\setlength{\parskip}{4pt}

% Title block: \wstitle{code}{title}{unit line}
\newcommand{\wstitle}[3]{\begin{center}{\small\bfseries\color{gold}WORKSHEET #1}\\[3pt]{\LARGE\bfseries\color{navy}#2}\\[3pt]{\small\color{black!60}#3}\end{center}\vspace{2pt}}
% Section heading
\newcommand{\sech}[1]{\par\vspace{10pt}{\large\bfseries\color{navy}#1}\par\nobreak\vspace{2pt}{\color{rule}\hrule height0.8pt}\vspace{7pt}}
% A level tag printed beside a group of questions
\newcommand{\level}[1]{\par\vspace{6pt}{\footnotesize\bfseries\color{gold}\MakeUppercase{#1}}\par\nobreak\vspace{2pt}}

\newtcolorbox{ideabox}{colframe=ink,colback=inkbg,coltitle=white,colbacktitle=ink,fonttitle=\bfseries\small,title={Key ideas},arc=2pt,boxrule=0.8pt,left=8pt,right=8pt}
\newtcolorbox{example}[1]{colframe=rule,colback=white,coltitle=navy,colbacktitle=inkbg,fonttitle=\bfseries\small,title={#1},arc=2pt,boxrule=0.8pt,borderline west={3pt}{0pt}{ink},left=10pt,right=8pt}
\newtcolorbox{frbox}[1]{colframe=gold,colback=goldbg,coltitle=white,colbacktitle=gold,fonttitle=\bfseries\small,title={#1},arc=2pt,boxrule=0.8pt,breakable,left=8pt,right=8pt}
\newtcolorbox{solbox}{colframe=okgreen,colback=okbg,coltitle=white,colbacktitle=okgreen,fonttitle=\bfseries\small,title={Worked solutions},arc=2pt,boxrule=0.8pt,breakable,left=8pt,right=8pt}
\newtcolorbox{scorebox}[1]{colframe=navy,colback=white,coltitle=white,colbacktitle=navy,fonttitle=\bfseries\small,title={#1},arc=2pt,boxrule=0.8pt,breakable,left=8pt,right=8pt}

\newcommand{\soln}{\par\smallskip\textbf{\color{ink}Solution.}\ }
\newcommand{\result}[1]{\par\smallskip\textbf{#1}}
% blank working space of a given height
\newcommand{\work}[1]{\par\vspace{#1}}
% a point value in the margin of a scoring guide
\newcommand{\pts}[1]{\hfill{\footnotesize\bfseries\color{gold}[#1]}}

% Coordinate grid for graph questions: \begin{apgraph}{xmin}{xmax}{ymin}{ymax} ... \end{apgraph}
\newenvironment{apgraph}[4]{\begin{center}\begin{tikzpicture}\begin{axis}[
  axis lines=middle, xmin=#1, xmax=#2, ymin=#3, ymax=#4,
  grid=both, grid style={black!12}, xtick distance=1, ytick distance=1,
  ticklabel style={font=\scriptsize}, xlabel={$x$}, ylabel={$y$},
  width=10cm, height=7.2cm, samples=80, clip=false,
  every axis plot/.append style={line width=1.1pt, color=ink}]}{\end{axis}\end{tikzpicture}\end{center}}
\newcommand{\closeddot}[2]{\addplot[only marks, mark=*, mark size=2.2pt, color=ink] coordinates {(#1,#2)};}
\newcommand{\opendot}[2]{\addplot[only marks, mark=*, mark size=2.2pt, color=ink, mark options={fill=white}] coordinates {(#1,#2)};}

% Multiple-choice options: four across, or stacked when the options are long.
\newcommand{\opts}[4]{\par\vspace{3pt}\begin{tabular}{@{}p{0.235\linewidth}p{0.235\linewidth}p{0.235\linewidth}p{0.235\linewidth}@{}}(A)\ #1 & (B)\ #2 & (C)\ #3 & (D)\ #4\end{tabular}\par\vspace{9pt}}
\newcommand{\optsv}[4]{\par\vspace{3pt}\begin{tabular}{@{}l@{\ \ }p{0.9\linewidth}@{}}(A) & #1\\(B) & #2\\(C) & #3\\(D) & #4\end{tabular}\par\vspace{9pt}}

\fancyhead[R]{\footnotesize\color{black!60}Worksheet 4.7}
\begin{document}

\wstitle{4.7}{Indeterminate Forms and L'H\^opital's Rule}{Unit 4: Derivatives in Context}

\begin{ideabox}
\begin{itemize}[leftmargin=*,itemsep=2pt,topsep=2pt]
  \item If $\dfrac{f(x)}{g(x)}$ has the form $\dfrac00$ or $\dfrac{\infty}{\infty}$, then $\displaystyle\lim\dfrac{f(x)}{g(x)}=\lim\dfrac{f'(x)}{g'(x)}$, provided the second limit exists.
  \item Differentiate numerator and denominator \textbf{separately}. This is not the quotient rule.
  \item \textbf{Check the form first}, and say so: state that both limits are 0 (or both infinite) before applying the rule.
  \item If the result is again indeterminate, apply the rule again.
\end{itemize}
\end{ideabox}

\sech{Worked examples}
\begin{example}{Example 1 \textperiodcentered{} The form 0/0}
Evaluate $\displaystyle\lim_{x\to 0}\dfrac{\sin(4x)}{3x}$.\soln Both $\sin(4x)$ and $3x$ tend to 0. By L'H\^opital's Rule the limit is $\displaystyle\lim_{x\to 0}\dfrac{4\cos(4x)}{3}$.\result{$\dfrac43$.}
\end{example}

\begin{example}{Example 2 \textperiodcentered{} The form infinity over infinity}
Evaluate $\displaystyle\lim_{x\to \infty}\dfrac{\ln x}{x}$.\soln Both tend to infinity. The rule gives $\displaystyle\lim_{x\to \infty}\dfrac{1/x}{1}$.\result{$0$.}
\end{example}

\begin{example}{Example 3 \textperiodcentered{} Applying the rule twice}
Evaluate $\displaystyle\lim_{x\to 0}\dfrac{e^{x}-1-x}{x^2}$.\soln Form $\tfrac00$: $\displaystyle\lim_{x\to 0}\dfrac{e^{x}-1}{2x}$, again $\tfrac00$: $\displaystyle\lim_{x\to 0}\dfrac{e^{x}}{2}$.\result{$\dfrac12$.}
\end{example}

\newpage
\sech{Practice}

\level{Core}
\begin{enumerate}[leftmargin=*,itemsep=2pt]
  \item Evaluate $\displaystyle\lim_{x\to 2}\dfrac{x^2-4}{x-2}$ using L'H\^opital's Rule.\work{2.3cm}
  \item Evaluate $\displaystyle\lim_{x\to 0}\dfrac{1-\cos x}{x}$.\work{2.3cm}
  \item Evaluate $\displaystyle\lim_{x\to 1}\dfrac{x^5-1}{x-1}$.\work{2.3cm}
  \item Can L'H\^opital's Rule be used for $\displaystyle\lim_{x\to 0}\dfrac{x+1}{x+3}$? Find the limit.\work{2.3cm}
\end{enumerate}

\level{Exam level}
\begin{enumerate}[leftmargin=*,itemsep=2pt,start=5]
  \item Evaluate $\displaystyle\lim_{x\to 0}\dfrac{\tan(3x)}{\sin(2x)}$.\work{2.8cm}
  \item Evaluate $\displaystyle\lim_{x\to \infty}\dfrac{x^2+1}{e^{x}}$.\work{2.8cm}
  \item Evaluate $\displaystyle\lim_{x\to 0}\dfrac{\sin x-x}{x^2}$.\work{2.8cm}
  \item A differentiable function $f$ has $f(2)=0$ and $f'(2)=3$, and $f'$ is continuous. Find $\displaystyle\lim_{x\to 2}\dfrac{f(x)}{x^3-8}$.\work{2.8cm}
\end{enumerate}

\begin{frbox}{Question 9 \textperiodcentered{} Free-response practice}
Evaluate each limit. In each part, show that L'H\^opital's Rule applies.
\begin{enumerate}[label=(\alph*),leftmargin=*,itemsep=2pt]
  \item $\displaystyle\lim_{x\to 0}\dfrac{e^{2x}-1}{\tan x}$\work{2.6cm}
  \item $\displaystyle\lim_{x\to \pi/2}\dfrac{\cos x}{x-\frac{\pi}{2}}$\work{2.6cm}
  \item $\displaystyle\lim_{x\to 0}\dfrac{g(x)}{e^{x}-1}$, where $g$ is differentiable with $g(0)=0$ and $g'(0)=5$, and $g'$ is continuous.\work{2.6cm}
\end{enumerate}
\end{frbox}

\newpage
\begin{solbox}
\begin{enumerate}[leftmargin=*,itemsep=4pt]
  \item Form $\tfrac00$: $\displaystyle\lim_{x\to 2}\dfrac{2x}{1}=4$.
  \item Form $\tfrac00$: $\displaystyle\lim_{x\to 0}\dfrac{\sin x}{1}=0$.
  \item Form $\tfrac00$: $\displaystyle\lim_{x\to 1}5x^4=5$.
  \item No: the numerator tends to 1 and the denominator to 3, so the form is not indeterminate. The limit is $\dfrac13$.
  \item Form $\tfrac00$: $\displaystyle\lim_{x\to 0}\dfrac{3\sec^2(3x)}{2\cos(2x)}=\dfrac32$.
  \item Form $\tfrac\infty\infty$, twice: $\displaystyle\lim_{x\to \infty}\dfrac{2x}{e^{x}}=\displaystyle\lim_{x\to \infty}\dfrac{2}{e^{x}}=0$.
  \item Form $\tfrac00$, twice: $\displaystyle\lim_{x\to 0}\dfrac{\cos x-1}{2x}=\displaystyle\lim_{x\to 0}\dfrac{-\sin x}{2}=0$.
  \item Both $f(x)$ and $x^3-8$ tend to 0. $\displaystyle\lim_{x\to 2}\dfrac{f'(x)}{3x^2}=\dfrac{3}{12}=\dfrac14$.
  \item \begin{enumerate}[label=(\alph*),leftmargin=*,itemsep=3pt]
    \item $e^{2x}-1\to0$ and $\tan x\to0$. $\displaystyle\lim_{x\to 0}\dfrac{2e^{2x}}{\sec^2x}=2$.
    \item $\cos x\to0$ and $x-\dfrac\pi2\to0$. $\displaystyle\lim_{x\to \pi/2}\dfrac{-\sin x}{1}=-1$.
    \item $g(x)\to g(0)=0$ because $g$ is continuous, and $e^{x}-1\to0$. $\displaystyle\lim_{x\to 0}\dfrac{g'(x)}{e^{x}}=\dfrac{5}{1}=5$.
  \end{enumerate}
\end{enumerate}
\end{solbox}

\end{document}
